% Part: normal-modal-logic
% Chapter: tableaux
% Section: more-soundness

\documentclass[../../../include/open-logic-section]{subfiles}

\begin{document}

\olfileid{nml}{tab}{msn}
      
\olsection{વધારાના નિયમોની યથાર્થતા}

નિયમને નિદર્શનોના કોઈ વર્ગ માટે યથાર્થ ત્યારે કહીએ જ્યારે
!!a{tableau}ની કોઈ શાખા તે વર્ગના નિદર્શનમાં સંતોષ્ય હોય ત્યારે
નિયમ લાગુ કરવાથી મળતી શાખા પણ તે વર્ગના નિદર્શનમાં સંતોષ્ય હોય.

\begin{prop}
  \ollabel{prop:soundness-T} 
  \iftag{prvBox}{\Ax{T}$\Box$}{}%
  \iftag{notprvBox,notprvDiamond}{}{ અને }%
  \iftag{prvDiamond}{\Ax{T}$\Diamond$}{}
  \iftag{notprvBox,notprvDiamond}{નિયમ}{નિયમો}
  સ્વવાચક નિદર્શનો માટે યથાર્થ છે.
\end{prop}

\begin{proof}
\begin{tagenumerate}{prvBox,prvDiamond}
\tagitem{prvBox}{$\sFmla{\True}{\Box !B}[\sigma] \in \Gamma$ને
  T$\Box$ લાગુ કરીને શાખા વિસ્તારીએ:
  \iftag{probBox}{અભ્યાસપ્રશ્ન.}{તેનાથી શાખા પર નવું !!{signed
      formula}~$\sFmla{\True}{!B}[\sigma]$ મળે છે. ધારો કે
    $\mSat{M}{\Gamma}[f]$, ખાસ કરીને $\mSat{M}{\Box
      !B}[f(\sigma)]$. $R$ સ્વવાચક હોવાથી
    $Rf(\sigma)f(\sigma)$. તેથી $\mSat{M}{!B}[f(\sigma)]$;
    એટલે $\mModel{M}, f$ એ $\sFmla{\True}{!B}[\sigma]$ને
    સંતોષે છે.}}{}
\tagitem{prvDiamond}{$\sFmla{\False}{\Diamond !B}[\sigma]
  \in \Gamma$ને T$\Diamond$ લાગુ કરીને શાખા વિસ્તારીએ:
  \iftag{probDiamond}{અભ્યાસપ્રશ્ન.}{તેનાથી શાખા પર નવું !!{signed
      formula}~$\sFmla{\False}{!B}[\sigma]$ મળે છે. ધારો કે
    $\mSat{M}{\Gamma}[f]$, ખાસ કરીને $\mSat/{M}{\Diamond
      !B}[f(\sigma)]$. $R$ સ્વવાચક હોવાથી
    $Rf(\sigma)f(\sigma)$. તેથી $\mSat/{M}{!B}[f(\sigma)]$;
    એટલે $\mModel{M}, f$ એ $\sFmla{\False}{!B}[\sigma]$ને
    સંતોષે છે.}}{}
\end{tagenumerate}
\end{proof}

\tagprob{probBox,probDiamond}
\begin{prob}
  \olref[nml][tab][msn]{prop:soundness-T}નું પ્રમાણ પૂર્ણ કરો.
\end{prob}
\tagendprob

\begin{prop}
  \ollabel{prop:soundness-D} 
  \iftag{prvBox}{\Ax{D}$\Box$}{}%
  \iftag{notprvBox,notprvDiamond}{}{ અને }%
  \iftag{prvDiamond}{\Ax{D}$\Diamond$}{}
  \iftag{notprvBox,notprvDiamond}{નિયમ}{નિયમો}
  સિરિયલ નિદર્શનો માટે યથાર્થ છે.
\end{prop}

\begin{proof}
\begin{tagenumerate}{prvBox,prvDiamond}
\tagitem{prvBox}{D$\Box$ને
  $\sFmla{\True}{\Box !B}[\sigma] \in \Gamma$ પર લાગુ કરીને
  શાખા વિસ્તારીએ:
  \iftag{probBox}{અભ્યાસપ્રશ્ન.}{તેનાથી શાખા પર નવું !!{signed
      formula}~$\sFmla{\True}{\Diamond !B}[\sigma]$ મળે છે.
    ધારો કે $\mSat{M}{\Gamma}[f]$, ખાસ કરીને
    $\mSat{M}{\Box !B}[f(\sigma)]$. $R$ સિરિયલ હોવાથી એવું
    $w \in W$ છે કે $Rf(\sigma)w$. ત્યારે $\mSat{M}{!B}[w]$;
    તેથી $\mSat{M}{\Diamond !B}[f(\sigma)]$.
    આમ $\mModel{M}, f$ એ
    $\sFmla{\True}{\Diamond!B}[\sigma]$ને સંતોષે છે.}}{}
\tagitem{prvDiamond}{D$\Diamond$ને
  $\sFmla{\False}{\Diamond !B}[\sigma] \in \Gamma$ પર લાગુ
  કરીને શાખા વિસ્તારીએ:
  \iftag{probDiamond}{અભ્યાસપ્રશ્ન.}{તેનાથી શાખા પર નવું !!{signed
      formula}~$\sFmla{\False}{\Box !B}[\sigma]$ મળે છે.
    ધારો કે $\mSat{M}{\Gamma}[f]$, ખાસ કરીને
    $\mSat/{M}{\Diamond !B}[f(\sigma)]$. $R$ સિરિયલ હોવાથી
    એવું $w \in W$ છે કે $Rf(\sigma)w$. ત્યારે
    $\mSat/{M}{!B}[w]$; તેથી
    $\mSat/{M}{\Box !B}[f(\sigma)]$. આમ $\mModel{M}, f$ એ
    $\sFmla{\False}{\Box!B}[\sigma]$ને સંતોષે છે.}}{}
\end{tagenumerate}
\end{proof}

\tagprob{probBox,probDiamond}
\begin{prob}
  \olref[nml][tab][msn]{prop:soundness-D}નું પ્રમાણ પૂર્ણ કરો.
\end{prob}
\tagendprob

\begin{prop}
  \ollabel{prop:soundness-B} 
  \iftag{prvBox}{\Ax{B}$\Box$}{}%
  \iftag{notprvBox,notprvDiamond}{}{ અને }%
  \iftag{prvDiamond}{\Ax{B}$\Diamond$}{}
  \iftag{notprvBox,notprvDiamond}{નિયમ}{નિયમો}
  સંમિત નિદર્શનો માટે યથાર્થ છે.
\end{prop}

\begin{proof}
\begin{tagenumerate}{prvBox,prvDiamond}
\tagitem{prvBox}{B$\Box$ને
  $\sFmla{\True}{\Box !B}[\sigma.n] \in \Gamma$ પર લાગુ કરીને
  શાખા વિસ્તારીએ:
  \iftag{probBox}{અભ્યાસપ્રશ્ન.}{તેનાથી શાખા પર નવું !!{signed
      formula}~$\sFmla{\True}{!B}[\sigma]$ મળે છે. ધારો કે
    $\mSat{M}{\Gamma}[f]$, ખાસ કરીને $\mSat{M}{\Box
      !B}[f(\sigma.n)]$. $f$ એ શાખાના ઉપસર્ગોનું
    $\mModel{M}$માં અર્થઘટન હોવાથી $Rf(\sigma)f(\sigma.n)$.
    $R$ સંમિત હોવાથી $Rf(\sigma.n)f(\sigma)$ પણ છે.
    $\mSat{M}{\Box !B}[f(\sigma.n)]$ હોવાથી
    $\mSat{M}{!B}[f(\sigma)]$. તેથી $\mModel{M}, f$ એ
    $\sFmla{\True}{!B}[\sigma]$ને સંતોષે છે.}}{}
\tagitem{prvDiamond}{B$\Diamond$ને
  $\sFmla{\False}{\Diamond !B}[\sigma.n] \in \Gamma$ પર લાગુ
  કરીને શાખા વિસ્તારીએ:
  \iftag{probDiamond}{અભ્યાસપ્રશ્ન.}{તેનાથી શાખા પર નવું !!{signed
      formula}~$\sFmla{\False}{!B}[\sigma]$ મળે છે. ધારો કે
    $\mSat{M}{\Gamma}[f]$, ખાસ કરીને $\mSat/{M}{\Diamond
      !B}[f(\sigma.n)]$. $f$ એ શાખાના ઉપસર્ગોનું
    $\mModel{M}$માં અર્થઘટન હોવાથી $Rf(\sigma)f(\sigma.n)$.
    $R$ સંમિત હોવાથી $Rf(\sigma.n)f(\sigma)$ પણ છે.
    $\mSat/{M}{\Diamond !B}[f(\sigma.n)]$ હોવાથી
    $\mSat/{M}{!B}[f(\sigma)]$. તેથી $\mModel{M}, f$ એ
    $\sFmla{\False}{!B}[\sigma]$ને સંતોષે છે.}}{}
\end{tagenumerate}
\end{proof}

\tagprob{probBox,probDiamond}
\begin{prob}
  \olref[nml][tab][msn]{prop:soundness-B}નું પ્રમાણ પૂર્ણ કરો.
\end{prob}
\tagendprob

\begin{prop}
  \ollabel{prop:soundness-4} 
  \iftag{prvBox}{\Ax{4}$\Box$}{}%
  \iftag{notprvBox,notprvDiamond}{}{ અને }%
  \iftag{prvDiamond}{\Ax{4}$\Diamond$}{}
  \iftag{notprvBox,notprvDiamond}{નિયમ}{નિયમો}
  પરંપરિત નિદર્શનો માટે યથાર્થ છે.
\end{prop}

\begin{proof}
\begin{tagenumerate}{prvBox,prvDiamond}
\tagitem{prvBox}{4$\Box$ને
  $\sFmla{\True}{\Box !B}[\sigma] \in \Gamma$ પર લાગુ કરીને
  શાખા વિસ્તારીએ:
  \iftag{probBox}{અભ્યાસપ્રશ્ન.}{તેનાથી શાખા પર નવું !!{signed
      formula}~$\sFmla{\True}{\Box!B}[\sigma.n]$ મળે છે.
    ધારો કે $\mSat{M}{\Gamma}[f]$, ખાસ કરીને
    $\mSat{M}{\Box !B}[f(\sigma)]$. $f$ એ શાખાના
    ઉપસર્ગોનું~$\mModel{M}$માં અર્થઘટન છે અને $\sigma.n$
    વપરાયેલો હોવો જોઈએ, તેથી $Rf(\sigma)f(\sigma.n)$.
    હવે $Rf(\sigma.n)w$ હોય એવું કોઈપણ જગત~$w$ લો.
    $R$ પરંપરિત હોવાથી $Rf(\sigma)w$.
    $\mSat{M}{\Box !B}[f(\sigma)]$ હોવાથી
    $\mSat{M}{!B}[w]$. તેથી
    $\mSat{M}{\Box !B}[f(\sigma.n)]$ અને $\mModel{M}, f$ એ
    $\sFmla{\True}{\Box !B}[\sigma.n]$ને સંતોષે છે.}}{}
\tagitem{prvDiamond}{4$\Diamond$ને
  $\sFmla{\False}{\Diamond !B}[\sigma] \in \Gamma$ પર લાગુ
  કરીને શાખા વિસ્તારીએ:
  \iftag{probDiamond}{અભ્યાસપ્રશ્ન.}{તેનાથી શાખા પર નવું !!{signed
      formula}~$\sFmla{\False}{\Diamond!B}[\sigma.n]$ મળે છે.
    ધારો કે $\mSat{M}{\Gamma}[f]$, ખાસ કરીને
    $\mSat/{M}{\Diamond !B}[f(\sigma)]$. $f$ એ શાખાના
    ઉપસર્ગોનું~$\mModel{M}$માં અર્થઘટન છે અને $\sigma.n$
    વપરાયેલો હોવો જોઈએ, તેથી $Rf(\sigma)f(\sigma.n)$.
    હવે $Rf(\sigma.n)w$ હોય એવું કોઈપણ જગત~$w$ લો.
    $R$ પરંપરિત હોવાથી $Rf(\sigma)w$.
    $\mSat/{M}{\Diamond !B}[f(\sigma)]$ હોવાથી
    $\mSat/{M}{!B}[w]$. તેથી
    $\mSat/{M}{\Diamond !B}[f(\sigma.n)]$ અને $\mModel{M}, f$ એ
    $\sFmla{\False}{\Diamond !B}[\sigma.n]$ને સંતોષે છે.}}{}
\end{tagenumerate}
\end{proof}

\tagprob{probBox,probDiamond}
\begin{prob}
  \olref[nml][tab][msn]{prop:soundness-4}નું પ્રમાણ પૂર્ણ કરો.
\end{prob}
\tagendprob

\begin{prop}
  \ollabel{prop:soundness-4r} 
  \iftag{prvBox}{\Ax{4r}$\Box$}{}%
  \iftag{notprvBox,notprvDiamond}{}{ અને }%
  \iftag{prvDiamond}{\Ax{4r}$\Diamond$}{}
  \iftag{notprvBox,notprvDiamond}{નિયમ}{નિયમો}
  યુક્લિડિયન નિદર્શનો માટે યથાર્થ છે.
\end{prop}

\begin{proof}
\sourcecorrection{OLMD-061}{મૂળમાં બંને યુક્લિડિયન કિસ્સામાં
$f(\sigma).n$ લખાયું છે, જે વિધેયના મૂલ્યને ઉપસર્ગ જેવું ગણે છે;
પૂર્વપક્ષ મુજબ અહીં $f(\sigma.n)$ જરૂરી છે.}
\begin{tagenumerate}{prvBox,prvDiamond}
\tagitem{prvBox}{4r$\Box$ને
  $\sFmla{\True}{\Box !B}[\sigma.n] \in \Gamma$ પર લાગુ કરીને
  શાખા વિસ્તારીએ:
  \iftag{probBox}{અભ્યાસપ્રશ્ન.}{તેનાથી શાખા પર નવું !!{signed
      formula}~$\sFmla{\True}{\Box!B}[\sigma]$ મળે છે. ધારો કે
    $\mSat{M}{\Gamma}[f]$, ખાસ કરીને $\mSat{M}{\Box
      !B}[f(\sigma.n)]$. $f$ એ શાખાના ઉપસર્ગોનું
    $\mModel{M}$માં અર્થઘટન હોવાથી $Rf(\sigma)f(\sigma.n)$.
    હવે $Rf(\sigma)w$ હોય એવું કોઈપણ જગત~$w$ લો.
    $R$ યુક્લિડિયન હોવાથી $Rf(\sigma.n)w$.
    $\mSat{M}{\Box !B}[f(\sigma.n)]$ હોવાથી
    $\mSat{M}{!B}[w]$. તેથી $\mSat{M}{\Box !B}[f(\sigma)]$;
    એટલે $\mModel{M}, f$ એ
    $\sFmla{\True}{\Box !B}[\sigma]$ને સંતોષે છે.}}{}
\tagitem{prvDiamond}{4r$\Diamond$ને
  $\sFmla{\False}{\Diamond !B}[\sigma.n] \in \Gamma$ પર લાગુ
  કરીને શાખા વિસ્તારીએ:
  \iftag{probDiamond}{અભ્યાસપ્રશ્ન.}{તેનાથી શાખા પર નવું !!{signed
      formula}~$\sFmla{\False}{\Diamond!B}[\sigma]$ મળે છે.
    \sourcecorrection{OLMD-062}{મૂળમાં ડાયમંડ-નિયમનો નિષ્કર્ષ
    $\sFmla{\True}{\Box!B}[\sigma]$ લખાયો છે; નિયમ-કોષ્ટક અને
    આ જ પ્રમાણના અંતિમ વાક્ય મુજબ તે
    $\sFmla{\False}{\Diamond!B}[\sigma]$ હોવો જોઈએ.}
    ધારો કે $\mSat{M}{\Gamma}[f]$, ખાસ કરીને
    $\mSat/{M}{\Diamond !B}[f(\sigma.n)]$.
    $f$ એ શાખાના ઉપસર્ગોનું $\mModel{M}$માં અર્થઘટન હોવાથી
    $Rf(\sigma)f(\sigma.n)$. હવે $Rf(\sigma)w$ હોય એવું
    કોઈપણ જગત~$w$ લો. $R$ યુક્લિડિયન હોવાથી
    $Rf(\sigma.n)w$. $\mSat/{M}{\Diamond !B}[f(\sigma.n)]$
    હોવાથી $\mSat/{M}{!B}[w]$. તેથી
    $\mSat/{M}{\Diamond !B}[f(\sigma)]$; એટલે
    $\mModel{M}, f$ એ
    $\sFmla{\False}{\Diamond !B}[\sigma]$ને સંતોષે છે.}}{}
\end{tagenumerate}
\end{proof}

\tagprob{probBox,probDiamond}
\begin{prob}
  \olref[nml][tab][msn]{prop:soundness-4r}નું પ્રમાણ પૂર્ણ કરો.
\end{prob}
\tagendprob

\begin{cor}
\ollabel{cor:soundness-logics} \olref[mru]{tab:logics-rules}માં
આપેલી !!{tableau} પ્રણાલીઓ નિદર્શનોના અનુરૂપ વર્ગો માટે
યથાર્થ છે.
\end{cor}

\end{document}
